\section{Regla de decisión}\label{sec:decision}
Con pocas parcelas, una idea puede parecer mejor sólo por azar (Sección~\ref{sec:problema}). Por eso fijamos la
regla antes de probar. Una idea se acepta sólo si cumple dos condiciones.
\begin{enumerate}
  \item \textbf{Le gana a la mejor versión anterior} en al menos 7 de cada 10 particiones de la validación cruzada,
        usando siempre las mismas particiones, estratificadas por municipio como la división oficial.
  \item \textbf{Pasa la prueba de permutación}, es decir, el modelo empeora cuando se revuelven sus variables nuevas.
\end{enumerate}
Las ideas que no cumplan la regla se registran como probadas y descartadas, porque también son un resultado.

\textbf{Orden de las pruebas.} Seguimos el orden de las hipótesis (Tabla~\ref{tab:orden}). La elección del algoritmo
y la combinación de modelos quedan al final, porque suelen dar las mejoras más pequeñas.

\begin{table}[!t]
\caption{Orden propuesto de las pruebas}\label{tab:orden}
\centering\footnotesize
\begin{tabularx}{\columnwidth}{@{}rX>{\raggedright\arraybackslash}p{2.1cm}@{}}
\toprule
Paso & Prueba & Mejora esperada \\
\midrule
0 & Fijar las particiones de la validación cruzada y la regla de decisión & Requisito \\
1 & Régimen como variable base (H1) & Grande \\
2 & Variables satelitales dentro de cada régimen (H3) & Pequeña a media \\
3 & Causa de las parcelas atípicas (H4) & Incierta, grande si aparece \\
4 & Entrenamiento tolerante al ruido (H2) & Pequeña \\
5 & Clima diario por parcela (H5) & Incierta \\
6 & Elección del algoritmo y combinación de modelos & Al final \\
\bottomrule
\end{tabularx}
\end{table}
